\chapter{Timing}
\label{ch:timing}

This chapter gives the timing inside the device: the clock, the sample period, the pattern in which the audio memory is
shared, and the deadlines that the pattern sets. Chapter~\ref{ch:signal-timing} gives the timing of the pins.

\section{Clocks and Periods}

The system oscillator of an NTSC system runs at 21,477,272 Hz. The device is clocked at half that rate by
\texttt{CLKI}, 10,738,636 Hz, a period of 93.1 ns. Everything the device does is a whole number of \texttt{CLKI}
cycles. The system oscillator of a PAL system runs at 21,281,370 Hz, so its \texttt{CLKI} is 10,640,685 Hz, a period of
94.0 ns. The figures in Table~\ref{tab:periods} are for an NTSC system; Section~\ref{sec:clock-rate} gives the PAL
values.

\begin{table}[H]
	\centering
	\begin{tabular}{p{3.6cm}p{3.6cm}p{5.2cm}}
		\toprule
		\textbf{Quantity}              & \textbf{Value}                               & \textbf{Derivation}                                  \\
		\midrule
		\texttt{CLKI} period           & 93.1 ns                                      & 2 oscillator periods                                 \\
		Output sample period           & 336 cycles, 31.29 $\mu$s                     & 672 oscillator periods                               \\
		Output sample rate             & 31,960.226 Hz                                & \texttt{CLKI} / 336                                  \\
		Envelope update interval       & 64 samples, 2.002 ms                         & 21,504 cycles; 499.4 updates a second for each voice \\
		Serial bit clock \texttt{BCLK} & 10 and 11 cycles alternately, 1.023 MHz mean & 32 bits in a sample period; 16 in each slot          \\
		Memory transaction             & 2 cycles, 186 ns                             & One read or one write                                \\
		Sequencer tick                 & \texttt{SEQ\_PERIOD} samples                 & Chosen by the program                                \\
		\bottomrule
	\end{tabular}
	\caption{Clocks and Periods}
	\label{tab:periods}
\end{table}

The output sample period is the device's hard deadline. Each period the device must read the next sample byte of every
voice that is playing a sample and produce the next stereo output sample. The sequencers have a softer one: their work
may run over several sample periods, but it must finish before the next tick (Section~\ref{sec:seq-deadlines}). The
sample generator and the memory schedule run continuously, whether the output is muted or not.

\section{Clock Rate}
\label{sec:clock-rate}

The device has one clock, \texttt{CLKI}, and every interval it measures is a count of \texttt{CLKI} cycles. It works at
any steady clock from 9.0 to 12.0~MHz. The system supplies the clock: a system whose oscillator differs from the one in
Table~\ref{tab:periods} divides it to a frequency in this range.

\begin{table}[H]
	\centering
	\begin{tabular}{lrrr}
		\toprule
		\textbf{Quantity}                     & \textbf{Minimum} & \textbf{Nominal} & \textbf{Maximum} \\
		\midrule
		\texttt{CLKI} frequency               & 9.000 MHz        & 10.738636 MHz    & 12.000 MHz       \\
		\texttt{CLKI} period                  & 111.1 ns         & 93.1 ns          & 83.3 ns          \\
		Output sample rate                    & 26,786 Hz        & 31,960 Hz        & 35,714 Hz        \\
		Envelope update interval              & 2.389 ms         & 2.002 ms         & 1.792 ms         \\
		Serial bit clock \texttt{BCLK} (mean) & 0.857 MHz        & 1.023 MHz        & 1.143 MHz        \\
		Memory transaction (2 cycles)         & 222.2 ns         & 186.2 ns         & 166.7 ns         \\
		Pitch against the nominal clock       & $-306$ cents     & 0                & $+192$ cents     \\
		\bottomrule
	\end{tabular}
	\caption{Clock Rate Range}
	\label{tab:clock-range}
\end{table}

Changing the clock changes every time the device keeps, in proportion: the sample rate, the envelope update interval,
the length of a tick of \texttt{SEQ\_PERIOD} samples, the rate of the serial output, and the pitch of every note. It
changes no behaviour measured in cycles or samples. The memory schedule, the number of instructions a sequence may run
in a tick, and the latency of a window access in cycles are the same at any clock in the range.

\subsection{Pitch at a Different Clock}

A pitch table (Section~\ref{sec:pitch-table}) gives each note a step in sample bytes per output sample, so the pitch
that a step produces is proportional to the sample rate, $\mathrm{CLKI} / 336$. A program that wants a note at an exact
frequency builds its tables for the clock of the system it runs on. A program that uses tables built for the nominal
clock plays every note sharp or flat in the ratio of the two clocks, by the amount in the last row of
Table~\ref{tab:clock-range} at the ends of the range. The tempo of a sequence changes by the same ratio. A program that
must keep time with the display, for example, counts display frames and not sequencer ticks.

In a PAL system the sample rate is 31,668.7~Hz, the ratio of the two clocks is 0.9909, and every note is 15.9 cents
flat of the same note in an NTSC system. Every time the device keeps, from an envelope update of 2.021~ms to a tick, is
0.9\% longer.

\subsection{What the Rest of the System Must Provide}

The device does not read or report its clock. The parts around it must keep up with the highest clock the system will
use:

\begin{itemize}
	\item \textbf{Audio memory.} The memory must complete an access within two \texttt{CLKI} cycles, less the set-up and
	      hold times of the interface, which are not yet specified. At the maximum clock this is 166.7 ns.
	\item \textbf{Converter.} The converter must accept a bit clock of 10 and 11 \texttt{CLKI} cycles alternately and a word rate of
	      \texttt{CLKI}/336 (Section~\ref{sec:converter-interface}).
	\item \textbf{Output filter.} The filter that follows the converter removes the images of the signal around the
	      sample rate. It is designed for the sample rate of the system, and must pass the highest frequency to be played
	      and stop the lowest image.
\end{itemize}

The clock must be continuous and steady while the device is running. The behaviour of the device when the clock
frequency changes while it runs is not yet specified.

\section{The Audio Memory Schedule}
\label{sec:memory-schedule}

The audio memory has two users, the sound processor and the CPU's window, and the device never makes them compete. It
divides time into a fixed pattern of eight cycles, made of four slots of two cycles. One slot belongs to the sound
processor and three to the CPU:

\begin{table}[H]
	\centering
	\begin{tabular}{lll}
		\toprule
		\textbf{Slot} & \textbf{Cycles of the pattern} & \textbf{Owner}  \\
		\midrule
		0             & 0--1                           & Sound processor \\
		1             & 2--3                           & CPU window      \\
		2             & 4--5                           & CPU window      \\
		3             & 6--7                           & CPU window      \\
		\bottomrule
	\end{tabular}
	\caption{The Audio Memory Schedule}
	\label{tab:memory-schedule}
\end{table}

\begin{figure}[H]
	\centering
	\begin{tikzpicture}[
			slot/.style={draw, minimum height=0.8cm, minimum width=3.0cm, align=center, font=\footnotesize},
			note/.style={font=\scriptsize, text=sircgray}
		]
		\foreach \i/\label/\fillc in {0/Sound processor/sirclightgray, 1/CPU window/white, 2/CPU window/white, 3/CPU window/white} {
				\node[slot, fill=\fillc] at (\i*3.0+1.5,0) {\label};
				\node[note] at (\i*3.0+1.5,-0.65) {cycles \number\numexpr\i*2\relax--\number\numexpr\i*2+1\relax};
			}
		\draw[thick] (0,-1.1) -- (12.0,-1.1);
		\node[note] at (6.0,-1.4) {one pattern: 8 cycles, repeated 42 times in each 336-cycle sample period};
	\end{tikzpicture}
	\caption{One Pattern of the Memory Schedule}
	\label{fig:memory-schedule}
\end{figure}

A sample period of 336 cycles holds exactly 42 patterns, so it holds 42 sound-processor slots and 126 CPU slots, in the
same places in every sample period. No slot can be borrowed: the sound processor never takes a CPU slot, and the CPU
never takes a sound-processor slot, whether or not the owner has anything to do. A slot that its owner does not use is
left idle.

\subsection{CPU Slots}
\label{sec:cpu-slots}

A CPU access to the audio memory is held until the next CPU slot, takes two cycles, and is then acknowledged. It is
acknowledged in the second to the fifth cycle of the access, counting the cycle in which the access is made as the
first (Table~\ref{tab:cpu-latency}). The time does not depend on the voices, the sequencers or the filter. The CPU
window can move 126 words, 252 bytes, in each sample period, which is 4.03 million words a second. How fast a block
move actually runs depends also on the bus cycle of the CPU (\hwmanual{}). Register accesses use no slot and are never
delayed.

\begin{table}[H]
	\centering
	\begin{tabular}{p{5.2cm}cccccccc}
		\toprule
		\textbf{Access made in cycle}  & 0 & 1 & 2 & 3 & 4 & 5 & 6 & 7 \\
		\midrule
		\textbf{Acknowledged in cycle} & 4 & 3 & 2 & 3 & 2 & 3 & 2 & 5 \\
		\bottomrule
	\end{tabular}
	\caption{CPU Window Access Time (Cycles of the Eight-Cycle Pattern)}
	\label{tab:cpu-latency}
\end{table}

An access made in the first cycle of a CPU slot (cycles 2, 4 and 6) is the quickest. The slowest is one made in the
last cycle of the pattern, which has to wait for the first CPU slot of the next.

\subsection{Sound Processor Slots}
\label{sec:spu-slots}

The sound processor uses its 42 slots for each sample period as follows, in this order of priority:

\begin{enumerate}
	\item \textbf{Sample reads.} A voice that is playing needs one word at most in each sample period, since it uses one
	      byte. Six voices need six slots at most. The reads are made first in the period.
	\item \textbf{Sequencer reads.} The slots that remain, at least 36 in each sample period, serve the sequencers. An
	      instruction takes one read. A note that loads a new instrument takes up to nine more: seven for the instrument
	      record and two for the pitch table entry; a note on the instrument the voice already has takes two.
\end{enumerate}

The voices share the sequencer slots in turn, so no voice waits for all of another's.

\section{Deadlines and Overrun}
\label{sec:deadlines}

The device sets the \texttt{MEM\_OVERRUN} bit of \texttt{STATUS} when it misses a deadline:

\begin{itemize}
	\item A sample read is not complete when the sample period ends. The slots set aside for sample reads make this impossible
	      while the device is working within its limits.
	\item The sequencer work started by a tick is not complete when the next tick comes (Section~\ref{sec:seq-deadlines}).
\end{itemize}

Work that takes longer than one sample period is not a missed deadline in itself. The bit is held until the program
clears it through \texttt{STATUSCLR}. An overrun never takes a slot from the CPU, and it never stops the voices: sample
reads go on, and a voice whose tick was missed is simply started again at the next one.

\subsection{Choosing a Tick}
\label{sec:choosing-tick}

How short a tick can be depends on how much the sequences do in each. A voice executes at most eight instructions in a
tick, and each instruction is one read; a note also reads its instrument record when the instrument has changed (seven
reads) and its pitch table entry (two). Six voices share at least 36 sequencer slots in each sample period, so the
number of reads in a tick, divided by 36, gives the shortest tick in sample periods that cannot overrun.
Table~\ref{tab:tick-limits} works out three cases.

\begin{table}[H]
	\centering
	\small
	\begin{tabular}{p{7.4cm}rrr}
		\toprule
		\textbf{Case, for each voice in one tick}                                                   & \textbf{Reads} & \textbf{Six voices} & \textbf{Periods at least} \\
		\midrule
		Worst: eight instructions, four of them notes on a different instrument                     & 44             & 264                 & 8                         \\
		A note on a new instrument and three other instructions                                     & 13             & 78                  & 3                         \\
		Notes on the instrument already loaded, with \texttt{INST}, \texttt{FINE} and \texttt{WAIT} & 6              & 36                  & 1                         \\
		\bottomrule
	\end{tabular}
	\caption{Reads in a Tick}
	\label{tab:tick-limits}
\end{table}

These are bounds. A sequence that waits between notes does less in most ticks, so a tick can be shorter than the table
says. A tick of 50 to 200 a second leaves nearly all of the slots idle.

\section{Latency}

A voice's sample byte is read during one sample period, scaled and mixed during the next, and sent to the converter
during the one after. The delay from a read to its sound at the output is three sample periods, 94 $\mu$s. A note takes
effect after the device has read its instrument record and pitch entry, which takes a sample period or two, and its
first sample byte then goes through the three periods above, so a note sounds within about five sample periods, 0.16
ms, of the sequencer reaching it. A note on a voice that already has its instrument is sooner. A write to a voice
register, such as volume, is seen from the next sample period.
